% Desarrollo separado, 9 de septiembre de 2026.
% Amplía el lector producto de centro_electronico_registros.tex.
% No identifica la media vuelta espacial con una conjugación física.
\subsubsection{Bidireccionalidad del lector producto y memoria de cocientes}
\label{bim09:seccion}

La lectura de producto conserva cuatro coordenadas aritméticas antes de
determinar una firma. Para $x=(i,j)\in D_9^2$, $D_9=\{1,\ldots,9\}$,
se definen
\[
 N(x)=ij,\qquad r(x)=1+((ij-1)\bmod9),\qquad
 q(x)=\frac{N(x)-r(x)}9,
\]
y el lector discreto
\[
 \ell(1,2,4,8,7,5)=(0,1,2,3,4,5),\qquad
 \ell(3)=\ell(6)=\ell(9)=0.
\]
Así, $N=r+9q$ conserva simultáneamente producto entero, residuo positivo
y cociente aritmético. La función $f_\ell(x)=\ell(r(x))$ es la
coordenada utilizada por la firma regional. Las funciones $N,r,q$ y
$f_\ell$ tienen dominio posicional; la dirección de transporte permanece
en la semilla orientada.

\paragraph{Avance, media vuelta y lectura anterior al transporte.}

Sea $\mathscr X_\times=D_9^2\times\{N,E,S,O\}$. Los vectores cardinales
son $\nu_N=(-1,0)$, $\nu_E=(0,1)$, $\nu_S=(1,0)$,
$\nu_O=(0,-1)$. La suma posicional se reduce a representantes positivos
módulo nueve. Definimos
\[
 P(x,d)=(x+\nu_d,d),\qquad H(x,d)=(x,\bar d),
 \qquad \nu_{\bar d}=-\nu_d.
\]
Estos operadores satisfacen $P^9=I$, $H^2=I$ y $HPH=P^{-1}$.

En las primeras cincuenta y cuatro transiciones TPK, el canal producto
se activa en las fases $2,5,8$ de cada ventana nonádica. Cada activación
lee la posición anterior al avance. Tras las transiciones $27$ y $54$
se aplica $H$. Por tanto, las dieciocho posiciones leídas son, para una
semilla $s$,
\[
 P^0s,P^1s,\ldots,P^8s,\quad
 HP^9s,\,P H P^9s,\ldots,P^8H P^9s,
\]
entendiendo que la lectura posicional omite únicamente la dirección.
Para cualquier función posicional $f$, pongamos $f_n=f(P^ns)$, con
índices módulo nueve. La palabra de lecturas crudas es
\begin{equation}
 \mathcal R_f(s)=
 (f_0,f_1,\ldots,f_8,\ f_0,f_8,f_7,\ldots,f_1).
 \label{bim09:palabra}
\end{equation}
Las seis sumas de ventana, cada una de tres activaciones, son
\begin{align}
 A_j^f(s)&=f_{3j}+f_{3j+1}+f_{3j+2},&&0\leq j<3,\label{bim09:ida}\\
 A_{3+j}^f(s)&=f_{-3j}+f_{-3j-1}+f_{-3j-2},&&0\leq j<3.
 \label{bim09:vuelta}
\end{align}
Para $f=f_\ell$, la firma es
$E_{\rm sig}(s)=([A_0^f]_{10},\ldots,[A_5^f]_{10})$.
El cociente decimal de cada acumulador es
$c_j^{10}=\lfloor A_j^f/10\rfloor$ y conserva
$A_j^f=[A_j^f]_{10}+10c_j^{10}$.

\begin{proposition}[Media vuelta y permutación de semiciclos]
\label{bim09:media-vuelta}
Para toda semilla orientada y toda función posicional $f$,
\[
 \mathcal R_f(Hs)=\operatorname{rot}_9\mathcal R_f(s),\qquad
 A_j^f(Hs)=A_{j+3\bmod6}^f(s).
\]
La misma permutación transporta los residuos y los cocientes decimales
de los acumuladores. Se aplica, en particular, a $N,r,q$ y $f_\ell$.
\end{proposition}
\begin{proof}
La relación $P^nH=HP^{-n}$ y la igualdad $f(Hs)=f(s)$ dan
$f(P^nHs)=f(P^{-n}s)=f_{-n}$. Sustituir en
\eqref{bim09:palabra} intercambia sus dos grupos de nueve entradas.
La suma en grupos de tres produce la identidad de ventanas. La división
euclídea por diez se efectúa componente a componente y respeta la
permutación.
\end{proof}

El avance también tiene una ley exacta sobre los acumuladores:
\begin{align}
 A_j^f(Ps)-A_j^f(s)&=f_{3j+3}-f_{3j},&&0\leq j<3,
 \label{bim09:avance-ida}\\
 A_{3+j}^f(Ps)-A_{3+j}^f(s)&=f_{1-3j}-f_{-2-3j},&&0\leq j<3.
 \label{bim09:avance-vuelta}
\end{align}
La prueba consiste en cancelar los dos sumandos comunes de cada ventana.
El transporte de una firma requiere, por tanto, las evaluaciones de
frontera que aparecen en estas diferencias, además de las sumas ya
expresadas.

\paragraph{Inversión cronológica y términos de frontera.}

Sea $\gamma=(x_0,\ldots,x_n)$ una trayectoria registrada y sea
$\gamma^\dagger=(x_n,\ldots,x_0)$ su inversión cronológica, con cada
arista orientada en sentido opuesto. Para una función posicional $f$,
la lectura anterior a cada avance es
\[
 C_f(\gamma)=\sum_{t=0}^{n-1}f(x_t).
\]
Entonces
\begin{equation}
 C_f(\gamma^\dagger)=C_f(\gamma)+f(x_n)-f(x_0).
 \label{bim09:frontera-inversion}
\end{equation}
Si $n=3m$ y $A_j^f$ agrupa tres aristas consecutivas,
\begin{equation}
 A_j^f(\gamma^\dagger)=A_{m-1-j}^f(\gamma)
 +f(x_{3(m-j)})-f(x_{3(m-j-1)}).
 \label{bim09:frontera-ventana}
\end{equation}
Ambas fórmulas se siguen de que la trayectoria inversa lee los extremos
de llegada de las aristas originales. Los sumandos interiores coinciden;
la diferencia conserva los extremos.

Para un observable de arista antisimétrico $a(x,y)=-a(y,x)$, su palabra
sí se transforma mediante
\[
 \mathcal R_a(\gamma^\dagger)
 =-\operatorname{rev}\mathcal R_a(\gamma).
\]
Esta identidad actúa sobre aristas orientadas. Las identidades
\eqref{bim09:frontera-inversion}--\eqref{bim09:frontera-ventana}
actúan sobre lecturas posicionales. La media vuelta $H$ cambia la
dirección inicial y conserva el orden del calendario; la inversión
cronológica invierte el registro completo. Sus dominios y sus reglas de
lectura quedan especificados por separado.

\paragraph{Cociclo de ocupación y memoria conservada en ida y vuelta.}

Para una arista cardinal $x\to y$, con representantes positivos
$\widetilde x,\widetilde y$, el cruce del recubrimiento es
\[
 b(x,y)=\frac{\widetilde x+\nu_d-\widetilde y}{9}\in\mathbb Z^2.
\]
Se cumple $b(y,x)=-b(x,y)$. En la categoría libre de caminos
registrados, definimos el lector aditivo
\begin{equation}
 \mathfrak m(\gamma)=
 \left(\sum_{t=0}^{n-1}b(x_t,x_{t+1}),\ C_q(\gamma),\ n\right).
 \label{bim09:memoria-aditiva}
\end{equation}
La concatenación satisface
$\mathfrak m(\gamma_2\circ\gamma_1)=
\mathfrak m(\gamma_1)+\mathfrak m(\gamma_2)$.
Es un cociclo aditivo sobre trayectorias registradas: el registro conserva
las aristas de ida y de vuelta como eventos diferentes.

Por \eqref{bim09:frontera-inversion},
\begin{equation}
 \mathfrak m(\gamma^\dagger\circ\gamma)=
 \bigl(0,\ 2C_q(\gamma)+q(x_n)-q(x_0),\ 2n\bigr),
 \label{bim09:memoria-retorno}
\end{equation}
donde $\gamma^\dagger\circ\gamma$ denota la concatenación temporal
de la ida seguida de la vuelta. La suma orientada de cruces se cancela;
la ocupación aritmética y la longitud permanecen. Este lector no
desciende al cociente que identifica un recorrido seguido de su inversa
con la trayectoria vacía, puesto que ese cociente elimina los eventos
que la segunda y la tercera coordenadas conservan.

Hay también una lectura par puramente fronteriza:
\[
 V_b(\gamma)=\sum_{t=0}^{n-1}
 b(x_t,x_{t+1})b(x_t,x_{t+1})^{\mathsf T},\qquad
 V_b(\gamma^\dagger\circ\gamma)=2V_b(\gamma).
\]
La lectura impar y la lectura par registran propiedades distintas de
las mismas aristas; ambas están calculadas antes de olvidar el camino.

\paragraph{Aplicación exacta a la fibra de firma \texorpdfstring{$555555$}{555555}.}

La enumeración finita de $\mathscr X_\times$ contiene $324$ semillas,
$26$ firmas y $24$ semillas con firma $555555$. La partición estable
por $P,H$ refina esta fibra en tres clases de ocho elementos. Pueden
etiquetarse por su firma después de un avance:
\[
 \mathscr C_{177},\quad\mathscr C_{717},\quad\mathscr C_{771}.
\]
La acción de media vuelta es
\[
 H\mathscr C_{177}=\mathscr C_{177},\qquad
 H\mathscr C_{717}=\mathscr C_{771},\qquad
 H\mathscr C_{771}=\mathscr C_{717}.
\]
El certificado finito adjunto verifica estas identidades sobre todos los
representantes. En las veinticuatro semillas, los seis acumuladores
crudos de $f_\ell$ valen exactamente cinco. Por tanto, los cocientes
decimales $c_j^{10}$ se anulan en esta fibra. La información adicional
se observa en las posiciones, las evaluaciones enteras y los cocientes
aritméticos, que tienen lecturas diferentes de esa firma.

La coordenada transversal fija $a$ del recorrido cardinal vale $4$ o
$5$ en esas semillas. Durante una vuelta de nueve avances se recorre
una vez cada valor de la coordenada móvil $u\in D_9$. Puesto que
$\gcd(a,9)=1$, los residuos positivos $r(au)$ permutan $D_9$.
De aquí se obtiene, sin utilizar una constante física,
\[
 \sum_{u=1}^9au=45a,\qquad
 \sum_{u=1}^9r(au)=45,\qquad
 \sum_{u=1}^9q(au)=5(a-1).
\]
La vuelta posterior recorre las mismas posiciones en sentido opuesto.
En las dieciocho activaciones,
\begin{equation}
 \sum N=90a,\qquad\sum r=90,\qquad
 \sum q=10(a-1)=
 \begin{cases}30,&a=4,\\40,&a=5.\end{cases}
 \label{bim09:ocupacion-30-40}
\end{equation}
El producto entero acumulado vale, respectivamente, $360$ o $450$.
El enrollamiento neto es cero; la suma de cocientes y las dieciocho
incidencias conservan la historia de la lectura. Cada una de las tres
clases de firma contiene representantes de ambos valores transversales.
La misma firma y la misma clase de transporte de ese cociente finito
pueden, por tanto, coexistir con ocupaciones aritméticas distintas.

Estos resultados conectan la memoria conservada con operaciones
concretas: permutación de semiciclos, diferencias de frontera,
cocientes aritméticos y registro de incidencias. La identificación
físico-matemática de cualquier lector posterior conserva sus mapas
propios; las identidades aquí demostradas pertenecen a la dinámica
interna del canal producto.
